\subsection{Coproducts}

In this section, X denotes a manifold of class \(C^{r}\) and x is a point of X.

13.5.1. Let \(k \leqslant r\) . If \(\Delta\) is the diagonal map \(y \mapsto (y, y)\) from X into \(X \times X\) , then \(\Delta_*\) maps \(\mathrm{T}_x^{(k)}(\mathrm{X})\) into \(\mathrm{T}_{(x,x)}^{(k)}(\mathrm{X} \times \mathrm{X})\) , which is a vector subspace of \(\mathrm{T}_{x}^{(k)}(\mathrm{X})\otimes\mathrm{T}_{x}^{(k)}(\mathrm{X})\) (cf. 13.4.5). We thus obtain a linear map, again denoted by \(\Delta_{*}\) (or \(c\)):

\[
\mathrm{T} _ {x} ^ {(k)} (\mathrm{X}) \rightarrow \mathrm{T} _ {x} ^ {(k)} (\mathrm{X}) \otimes \mathrm{T} _ {x} ^ {(k)} (\mathrm{X});
\]

we call it the coproduct attached to \(\mathbf{T}_{x}^{(k)}(\mathbf{X})\). Equipped with this coproduct, \(\mathbf{T}_{x}^{(k)}(\mathbf{X})\) is a coassociative and cocommutative coalgebra (A, III, p. 144–145); it admits as its counit the linear form which assigns to a point distribution its constant term (13.2.1). If \(k' \leqslant k\) , the inclusion of \(\mathbf{T}_{x}^{(k')}(\mathbf{X})\) into \(\mathbf{T}_{x}^{(k)}(\mathbf{X})\) is a morphism of coalgebras.

If \(c = (\mathrm{U},\varphi ,\mathrm{E})\) is a chart of \(\mathbf{X}\) centered at \(x\) , the isomorphism

\[
\theta_ {c} \colon \mathrm{T} ^ {(k)} (\mathrm{E}) \rightarrow \mathrm{T} _ {x} ^ {(k)} (\mathrm{X})
\]

is an isomorphism of coalgebras, \(\mathrm{T}^{(k)}(\mathrm{E})\) being equipped with the coproduct induced by that of TS(E), cf. A, IV, § 5, no. 7.

If \(\varphi: X \to Y\) is a morphism of manifolds of class \(C^r\) , the map from \(T_x^{(k)}(X)\) to \(T_{\varphi(x)}^{(k)}(Y)\) induced by \(\varphi_*\) is a morphism of coalgebras.

13.5.2. Let \(F_{1}\) , \(F_{2}\) and F be separated polynormed spaces, and let \((u_{1}, u_{2}) \mapsto u_{1}.u_{2}\) be a continuous bilinear map from \(F_{1} \times F_{2}\) to F. Let \(t \in \mathrm{T}_{x}^{(r)}(\mathrm{X})\) and let \(\boldsymbol{c}(t) = \sum_{j} u_{j} \otimes v_{j}\), where \(u_{j},v_{j} \in \mathrm{T}_{x}^{(r)}(\mathrm{X})\), be its image under the coproduct. Let \(f_{i}: X \to F_{i}\) (\(i=1,2\)) be maps of class \(C^{r}\) . One has

\[
\langle f _ {1}. f _ {2}, t \rangle = \sum_ {j} \langle f _ {1}, u _ {j} \rangle . \langle f _ {2}, v _ {j} \rangle ,
\]

which we write, by abuse of notation:

\[
\langle f _ {1}. f _ {2}, t \rangle = \langle f _ {1} \otimes f _ {2}, c (t) \rangle .
\]

13.5.3. Let \(t \in \mathbf{T}_x^{(k)}(\mathbf{X})\) , where \(k \leqslant r\) .

\begin{enumerate}
\def\labelenumi{\alph{enumi})}
\item
  For \(c(t)=t\otimes t\) it is necessary and sufficient that \(t\) be equal to 0 or to \(\varepsilon_{x}\) .
\item
  For \(c(t)=t\otimes\varepsilon_{x}+\varepsilon_{x}\otimes t\) it is necessary and sufficient that \(t\) be a tangent vector, i.e., that one has \(t\in\mathrm{T}_{x}^{(1)+}(X)\) .
\end{enumerate}

